\documentclass[11pt]{amsart}
\usepackage{amsmath,amssymb,amsthm}
\usepackage{booktabs}
\usepackage{enumitem}
\usepackage[margin=0.9in]{geometry}
\usepackage[colorlinks=true,linkcolor=blue,citecolor=blue,urlcolor=blue]{hyperref}
\setlength{\emergencystretch}{3em}
\raggedbottom
\pdfstringdefDisableCommands{\def\e{eps}\def\vartheta{theta}}

\newtheorem{theorem}{Theorem}[section]
\newtheorem{lemma}[theorem]{Lemma}
\newtheorem{proposition}[theorem]{Proposition}
\newtheorem{corollary}[theorem]{Corollary}
\theoremstyle{definition}
\newtheorem{definition}[theorem]{Definition}
\theoremstyle{remark}
\newtheorem{remark}[theorem]{Remark}

\newcommand{\e}{\varepsilon}
\newcommand{\cX}{\mathcal{X}}
\newcommand{\cQ}{\mathcal{Q}}
\newcommand{\Lfp}{Lemma~\ref{lem:L44}}

\title[M\"obius in all intervals of length $x^{0.5435}$]{M\"obius cancellation in all intervals of length $x^{0.5435}$}
\author{Automated multi-agent proof search (Claude, Anthropic)}
\date{3 October 2026 (phase~II version). Refereed internally by independent agents (R2, R3; phase~II changes: R4), not by humans. See \S\ref{sec:prov} for provenance and status.}



\begin{document}
\begin{abstract}
We give an argument, partly computer-assisted and refereed only internally, that for every $\vartheta>0.5435$ and $\e_0>0$ one has
$\sum_{x<n\le x+H}\mu(n)\ll H(\log x)^{-1/3+\e_0}$ for all $H\ge x^{\vartheta}$. This would lower the exponent $0.55$
of Matom\"aki and Ter\"av\"ainen by $0.0065$. Their reduction is kept unchanged. The only new step replaces their Heath-Brown--Iwaniec type lemma.
In its place we use a level-set argument built on the Guth--Maynard large values estimates, including the $\inf_k$ form of their
Proposition~12.1, and on the fourth moment of $\zeta$. A branch-and-bound certificate in exact rational arithmetic, with margin $10^{-5}$,
closes the exponent bookkeeping. The binding configuration is five smooth factors of length $x^{1/5}$. With our inputs it closes exactly when
$\vartheta>69/127=0.5433071\ldots$, which is the limit of the method. The certificate also passes at $\vartheta=0.543308=69/127+9.1\cdot10^{-7}$,
and a control run at $\vartheta=0.5433$ fails. We also record why the routes tried in this project do not go below $69/127$; in the binding case
the problem is a fifth-moment bound for a single segment of the zeta function.
\end{abstract}
\maketitle

\section{Introduction}

Let $\mu$ be the M\"obius function. Matom\"aki and Ter\"av\"ainen \cite[Theorem~1.1]{MT} proved the following. \emph{Let $\theta>0.55$ and
$\e>0$ be fixed. Then for $x$ large enough and $H\ge x^{\theta}$,
$\sum_{x<n\le x+H}\mu(n)=O(H/(\log x)^{1/3-\e})$.} The exponent $0.55=11/20$ is the Heath-Brown--Iwaniec exponent for primes
in short intervals \cite{HBI}. It enters \cite{MT} through one lemma, \cite[Lemma~4.4]{MT}, which they take from
\cite[Lemma~2]{HBI} (see also \cite[Lemma~10.12]{Ha}). MT remark in \cite[\S2.1]{MT} that the proof of that lemma is not continuous in
$\theta$ and depends essentially on $\theta>0.55$. This is a paraphrase of text obtained through a summarising fetch tool.
In the level-set model used below, the classical large values inputs (the mean value theorem, Huxley's estimate \cite{Hu} and
Jutila's estimate \cite{Ju}) are consistent with this exponent: the same code certifies $0.552$ and fails at $0.548$ (\S\ref{sec:limits}). Guth and Maynard \cite{GM} have since improved large
values estimates for Dirichlet polynomials taking values near $N^{3/4}$. This note records what their estimates give in the Matom\"aki--Ter\"av\"ainen
framework. An earlier version of this note proved the result below for $\vartheta>0.548$, using only the closed form of \cite[Proposition~12.1]{GM}.

\begin{theorem}\label{thm:main}
Let $\vartheta>0.5435$ and $\e_0>0$ be fixed. Then for $x$ large enough and $H\ge x^{\vartheta}$,
\[
\sum_{x<n\le x+H}\mu(n)\ll \frac{H}{(\log x)^{1/3-\e_0}}.
\]
In particular, the exponent $11/20$ of \cite{MT} can be lowered by $0.0065$: for every fixed $\vartheta>11/20-0.0065$ one has
$\sum_{x<n\le x+H}\mu(n)=o(H)$ for all $H\ge x^{\vartheta}$.
\end{theorem}

\noindent\textbf{Status.} The proof (\S\S\ref{sec:red}--\ref{sec:proof}) is complete as written, modulo three things.
\begin{enumerate}[label=(\alph*),leftmargin=*]
\item The cited results. The statements from \cite{MT} were read only through a summarising tool. An independent re-check (S3), also via
such a tool, agrees (Remark~\ref{rem:sources}). Classical results are cited as standard, without theorem numbers.
\item A computer certificate (\S\ref{sec:cert}) in exact rational arithmetic, with margin $10^{-5}$. It extends the certificate of the
$0.548$ version (run in double precision and ported to exact arithmetic by a separate agent) by one estimate, the $\inf_k$ form of
\cite[Proposition~12.1]{GM}. The extension was written by one agent (GA). The referee R4 re-ran it, and re-checked every accepted
leaf with an independently written script (\S\ref{sec:cert}); the writer of this note also re-ran the headline invocation and obtained the same tiling.
\item Refereeing. The proofs of the $0.548$ version were refereed by an independent agent (R2), with verdict \emph{correct with minor fixable
issues}; all issues were addressed. A second, fresh-eyes referee (R3) read that note; its minor issues were also applied. The changes made for
$0.5435$ (listed in \S\ref{sec:prov}) were refereed by a third independent agent (R4), with verdict \emph{accept}: no blocking or major
issues, and four minor ones, all addressed here. There has been no human refereeing.
\end{enumerate}
We regard the result as \emph{unverified} until it has been checked outside this project.

\smallskip
\noindent\textbf{Method.} We follow \cite{MT} verbatim up to the point where \cite[Lemma~4.4]{MT} is used (\S\ref{sec:red}). In its place we prove
\Lfp\ (\S\ref{sec:L44}). The configuration is a product $ABC$ of Dirichlet polynomials: $A$ and $B$ are arbitrary and divisor-bounded,
and $C$ is a smooth (zeta-type) partial sum. We need an $L^1$ bound for $ABC$ on $[N_C,x^{\tau}]$ with $\tau\approx 0.4565$.
We discretise and split into level sets. On each level set the number of points is bounded by the smallest of several estimates:
\begin{itemize}[leftmargin=*]
\item the mean value theorem;
\item Guth--Maynard's Theorem~1.1 and the $\inf_k$ form of their Proposition~12.1, applied to a fixed finite family of products of $A$, $B$, $C$;
\item a fourth-moment bound for $C$, which we transfer from $\zeta$ (Lemma~\ref{lem:Z}).
\end{itemize}
Three regions need separate arguments. Small $t$ and long $C$ are handled by Cauchy--Schwarz. The ``corner'' where $|A|$ and $|B|$ are both nearly maximal is
handled by the Hal\'asz--Montgomery inequality and the Vinogradov--Korobov saving of the prime polynomial. A rigorous box
certificate covers the rest of parameter space.

\smallskip
\noindent\textbf{Limits.} With smaller margins the same certificate passes at $\vartheta=0.543308=69/127+9.1\cdot10^{-7}$
(Remark~\ref{rem:543308}). The number $69/127=0.5433071\ldots$ is the limit of the method. For five smooth factors of length $x^{1/5}$ the
inputs used close the case exactly when $\vartheta>69/127$; this is a refereed computation (P2, R1), recomputed exactly in phase~II (PB1).
A control run of the certificate at $\vartheta=0.5433$ fails. Section~\ref{sec:limits} records what phase~II of the search found about going
below $69/127$. In the binding case the problem is a fifth-moment bound for one zeta segment. A joint large values route is refuted. An
additive-energy input inside the Guth--Maynard method could gain at most $69/127-19/35\approx0.00045$. We also list priced targets.
Finally, \S\ref{sec:sub542} records a refereed reduction below $69/127$. For $\vartheta>0.542$, everything outside a near-equal family of
configurations is certified.

\section{The Matom\"aki--Ter\"av\"ainen reduction}\label{sec:red}

Throughout, $T_0=\exp((\log x)^{1/3})$. Write $f\ll_A g$ when the constant depends on $A$; all other constants may depend on the fixed
parameters. We use MT's notation with a general exponent. Fix $\theta\in(0.53,0.56)$ and a small $\e>0$. Let $H\asymp
x^{\theta+\e}$. Take $P=\exp((\log x)^{2/3+\e/2})$, $Q=x^{1/(\log\log x)^2}$ and $k=20$, as in \cite[\S4.1]{MT}. MT take $\theta=0.55$.

\begin{lemma}[{\cite[Lemma~3.1]{MT}}]\label{lem:ramare}
Let $\e>0$ be fixed, $x$ large, $(\log x)^4\le P<Q\le x^{o(1/\log\log x)}$ and $x^{\e}\le H\le x$. Then
\[
\sum_{x<n\le x+H}\mu(n)=-\!\!\sum_{\substack{x<prn\le x+H\\ P<p\le Q,\ r\le x^{\e/2}}}\!\! a_r\mu(n)+O\Big(H\frac{\log P}{\log Q}\Big),
\]
with explicit divisor-bounded coefficients $a_r$ (a Ramar\'e-type weight).
\end{lemma}

With the above $P,Q$, the error term is $\ll H(\log x)^{-1/3+\e}$. This is the only source of the exponent $1/3$.
Heath-Brown's identity with $k=20$ and a dyadic decomposition \cite[\S4.1]{MT} reduce matters to $\ll(\log x)^{2k+2}$ sums over
$x<prn_1\cdots n_{2k-1}\le x+H$, with $p\sim P_1$, $r\sim R$ and $n_i\sim N_i$. The parameters satisfy
\begin{equation}\label{eq:config}
P_1\in[P,Q],\quad R\in[\tfrac12,x^{\e/2}],\quad N_i\in[\tfrac12,x]\ (i\le k-1),\quad N_i\in[\tfrac12,(2x)^{1/k}]\ (i\ge k),\quad P_1R\textstyle\prod_iN_i\asymp x.
\end{equation}
For $i\le k-1$ the coefficients $a_i(n)$ are $1$ or $\log n$; we call these factors \emph{smooth}. For $i\ge k$ they are of M\"obius type,
$\mu(n)\mathbf 1_{n\le(2x)^{1/k}}$. Each such sum must be compared with $H/y_1$ times the same sum over $(x,x+y_1]$, where
$y_1=x\exp(-3(\log x)^{1/3})$; call this comparison $(\ast)$. Write
$P(s)=\sum_{P_1<p\le 2P_1}p^{-s}$, $R(s)=\sum_{r\sim R}a_rr^{-s}$ and $N_i(s)=\sum_{n\sim N_i}a_i(n)n^{-s}$.

\begin{lemma}[{\cite[Lemma~4.2]{MT}}]\label{lem:perron}
Suppose that for all $A>0$,
\begin{equation}\label{eq:L1}
\int_{T_0}^{x^{1+\e/10}/H}|P R N_1\cdots N_{2k-1}(\tfrac12+it)|\,dt\ll_A x^{1/2}(\log x)^{-A}.
\end{equation}
Then $(\ast)$ holds with error $O_A(H(\log x)^{-A})$.
\end{lemma}

\begin{lemma}[Vinogradov--Korobov; proof of {\cite[Lemma~4.3]{MT}}]\label{lem:VK}
For $|t|\le x^2$, $|P(1+it)|\ll P_1^{-1/(\log x)^{2/3+\e/3}}(\log x)^3+\log x/(|t|+1)$. The same bound holds for every
partial sum over $(P_1,u]$ with $u\le2P_1$. Partial summation then gives $\sup_{T_0\le|t|\le x^2}|P(\frac12+it)|\ll_A P_1^{1/2}(\log x)^{-A}$.
This is the step MT use in their proof of \cite[Lemma~4.3]{MT}, whose final bound carries the factor $P_1^{1/2}(\log x)^{-A}$.
\end{lemma}

\begin{lemma}[Type II; {\cite[Lemma~4.3]{MT}} with a general window]\label{lem:typeII}
Let $T=x^{1+\e/10}/H$ and $X_I=\prod_{i\in I}N_i$. If some $I\subset[2k-1]$ has both $X_I\ge T$ and $X_{[2k-1]\setminus I}\ge T$, then
\eqref{eq:L1} holds.
\end{lemma}

\begin{proof}
This is MT's proof. Lemma~\ref{lem:VK} bounds $P$ and the trivial bound handles $R$. The rest is
Cauchy--Schwarz and the mean value theorem, Lemma~\ref{lem:MV}:
$\int_{T_0}^T|\prod_I N_i\prod_{I^c}N_i|\ll(T+X_I)^{1/2}(T+X_{I^c})^{1/2}(\log x)^{O(1)}\ll(X_IX_{I^c})^{1/2}(\log x)^{O(1)}$.
\end{proof}

MT state the window as $X_I\in[x^{0.45-\e/2},x^{0.55+\e/2}]$. At its upper end the complement condition holds only after a
harmless adjustment: either group $R$ with $I^c$, or note that MT's Lemma~4.5 only produces $X_I\in[x^{0.45},x^{0.55}]$.

\begin{lemma}[Combinatorics; cf.\ proof of {\cite[Lemma~4.5]{MT}}]\label{lem:comb}
Let $0\le w_1<w_2$, $\alpha_i>0$ and $S=\sum_i\alpha_i>w_2$. Suppose no subsum lies in $[w_1,w_2]$. Let $I$ have the largest subsum
$\le w_2$. Then $\sum_I\alpha_i<w_1$. Moreover, some $r\notin I$ exists, and every $r\notin I$ satisfies $\alpha_r>w_2-\sum_I\alpha_i$ and
$\sum_{J}\alpha_j<S-w_2$, where $J=[K]\setminus(I\cup\{r\})$.
\end{lemma}

\begin{proof}
$\sum_I\le w_2$ is not in $[w_1,w_2]$, so it is $<w_1$. Since $\sum_I<S$, some $r\notin I$ exists. Maximality and
$\alpha_r>0$ give $\sum_I+\alpha_r>w_2$, and then $\sum_J=S-\sum_I-\alpha_r<S-w_2$.
\end{proof}

\begin{remark}[$\theta$-independence of the rest of MT's proof; audit]\label{prop:indep}
In the proof of \cite[Theorem~1.1]{MT}, the hypothesis $\theta>0.55$ enters only in five places (lemma numbers refer to \cite{MT}):
\begin{enumerate}[label=(\roman*),leftmargin=*]
\item Lemma~4.4;
\item the window of Lemma~4.3, which comes from the cut $T=x^{1+\e/10}/H$;
\item the numbers in Lemma~4.5;
\item the type~I Lemma~4.1, which MT note is already covered by the other two lemmas (paraphrase);
\item the check that $k=20$ and $N_r\gg x^{0.079}$ in the proof of Lemma~4.4.
\end{enumerate}
Everything else uses only $x^{\e}\le H\le x$:
\begin{itemize}[leftmargin=*]
\item Lemma~\ref{lem:ramare};
\item Heath-Brown's identity and the reduction to $H\asymp x^{\theta+\e}$;
\item Lemma~\ref{lem:perron}, including the ranges $|t|\le T_0$ and the polynomial $R$;
\item Lemma~\ref{lem:VK};
\item the final step through the prime number theorem with Vinogradov--Korobov error term.
\end{itemize}
\end{remark}

This is the conclusion of an audit of \cite{MT} (internal report S2). An independent re-check of the MT text (S3) agrees that only
\cite[Lemma~4.4]{MT} needs $\theta>0.55$; its quotes were also obtained through a summarising web tool. The S2 audit used the arXiv PDF and the ar5iv rendering,
read through a summarising fetch tool. For Lemma~\ref{lem:perron}, MT cite \cite[Lemma~7.2]{Ha}, which we could not consult. The audit
instead re-derived the Perron step and found it uniform in $\theta$. So, by Remark~\ref{prop:indep}, items (i)--(v) are the only
changes needed. Lemma~\ref{lem:typeII} and Lemma~\ref{lem:comb} take care of (ii)--(iv). Item (v) holds because $2\theta-1-\e>1/20$ for
$\theta>0.53$. Item (i) is \Lfp\ below.

\section{Analytic inputs}\label{sec:inputs}

For each input we give the statement in the form used. ``$\lessapprox$'' means $\le C_\eta T^{\eta}$ for every $\eta>0$.

\begin{lemma}[Mean value theorem; standard, see {\cite{MV}}, {\cite{Ga}}]\label{lem:MV}
$\int_0^U|\sum_{n\le N}a_nn^{-it}|^2dt=\sum|a_n|^2(U+O(n))$. Also, if the $t_r\in[0,T]$ are $1$-separated, then
$\sum_r|\sum_{n\le N}a_nn^{-it_r}|^2\ll(T+N)\log(2N)\sum|a_n|^2$.
\end{lemma}

\begin{lemma}[{Guth--Maynard \cite[Theorem~1.1]{GM}}]\label{lem:GM1}
Suppose $|b_n|\le1$ and $(t_r)_{r\le R}$ are $1$-separated points in $[0,T]$ with $|\sum_{n=N}^{2N}b_nn^{it_r}|\ge V$ for all $r\le R$. Then
$R\le T^{o(1)}(N^2V^{-2}+N^{18/5}V^{-4}+TN^{12/5}V^{-4})$.
\end{lemma}

\begin{lemma}[{Guth--Maynard \cite[Proposition~12.1]{GM}}, quoted verbatim]\label{lem:GM12}
``Suppose $(b_n)_{n\sim N}$, $(t_r)_{r\le R}$ are as in Theorem~1.1, and that $T^{5/6}\le N\le T$ and $V=N^{\sigma}$ with $\sigma\ge7/10$. Then
we have
\[
R\lessapprox N^{2-2\sigma}+T^{1/2}N^{3-4\sigma}+\inf_{k\in\mathbb N}\Big(T^{\frac{k}{k+1}}N^{(4-6\sigma)\frac{k}{k+1}}+N^{(5-6\sigma)\frac{4k}{4k+3}}T^{\frac{2}{4k+3}}\Big).
\]
In particular, we have
\[
R\lessapprox N^{2-2\sigma}+T^{1/2}N^{3-4\sigma}+T^{(30\sigma-21)/5}N^{(46-60\sigma)/5}."
\]
\end{lemma}

Here ``Theorem~1.1'' is \cite[Theorem~1.1]{GM}, i.e.\ Lemma~\ref{lem:GM1}: the hypotheses are $|b_n|\le1$, the $t_r$ are $1$-separated in
$[0,T]$, and $|\sum_{n=N}^{2N}b_nn^{it_r}|\ge V$ for every $r$. In \cite[\S1.2]{GM}, $A\lessapprox B$ means that for every $\epsilon>0$ there is
$C(\epsilon)$ with $|A|\le C(\epsilon)T^{\epsilon}B$ for all large $T$. The quotation was compared line by line with the arXiv text of \cite{GM}
(reports GA and R4; the closed form was also checked by R2 for the $0.548$ version). We use the first display with $k\in\{1,\dots,12\}$ only. Since the
infimum is at most its value at any fixed $k$, each fixed $k$ gives a valid bound, with a $\lessapprox$-constant independent of $k$ (R4).
The proof does not use the ``in particular'' display (the closed form used in the $0.548$ version); it enters only one cross-check run (\S\ref{sec:cert}).

\begin{lemma}[Ingham {\cite{In}}; standard, see also {\cite{Ti}}]\label{lem:Ing}
$\int_0^T|\zeta(\frac12+it)|^4dt\ll T\log^4T$.
\end{lemma}

\begin{lemma}[Hal\'asz--Montgomery inequality; standard, see {\cite{Mo}}, {\cite{Iv}}]\label{lem:HM}
If the $t_r\in[-T,T]$ are $1$-separated and $c_n$ is supported on $n\asymp N$, then
$\sum_{r\le R}|\sum_nc_nn^{-it_r}|^2\ll(N+RT^{1/2})\log(2T)\sum|c_n|^2$.
\end{lemma}

\begin{lemma}[Approximation by partial sums; standard, see {\cite{Ti}}]\label{lem:T411}
For $\sigma\ge\sigma_0>0$, $|t|<2\pi y/C$ with $C>1$ fixed: $\zeta(s)=\sum_{n\le y}n^{-s}-\frac{y^{1-s}}{1-s}+O(y^{-\sigma})$.
\end{lemma}

\begin{lemma}[Third-derivative test; standard, see {\cite{Ti}}, {\cite{GK}}]\label{lem:T3}
Let $f$ be real with $\lambda_3\le|f'''|\le h\lambda_3$ on $[a,b]$. Then $\sum_{a<n\le b}e(f(n))\ll h^{1/2}(b-a)\lambda_3^{1/6}+(b-a)^{1/2}\lambda_3^{-1/6}$.
\end{lemma}

We also use two standard facts, see \cite{Ti}: the truncated Perron formula and the convexity bound $\zeta(\sigma+it)\ll|t|^{(1-\sigma)/2+o(1)}$
for $\frac12\le\sigma\le1$.

\begin{remark}\label{rem:sources}
We could not re-check theorem numbers in \cite{Ti}, \cite{Mo}, \cite{Iv} and \cite{GK} while writing. So the classical results are cited
as standard, without numbers. The referee R2 confirmed that their statements are standard and correct. Lemmas~\ref{lem:GM1} and~\ref{lem:GM12}
were checked against the arXiv text of \cite{GM}: both by R2 and the earlier writer for the $0.548$ version, and the full $\inf_k$ display again in
phase~II (GA, and this writer). The statements from \cite{MT} were checked only through a
summarising web tool (audit S2). An independent re-check (S3), also through such a tool, agrees.
\end{remark}

\section{The replacement for the Heath-Brown--Iwaniec lemma}\label{sec:L44}

\subsection{Statement}
Let $T=x^{\tau}$. Let $A(s)=\sum\alpha_nn^{-s}$ with $\alpha_n$ supported in $(N_A,2^{40}N_A]$ and $|\alpha_n|\le d_K(n)(\log 2n)^K$.
Define $B$ in the same way with $\beta_m$ and $N_B$. Let $C(s)=\sum_{N_C<n\le2N_C}c_nn^{-s}$ with $c_n\equiv1$ or $c_n=\log n$. Write $N_A=x^a$,
$N_B=x^b$, $N_C=x^c$ and $s_0=a+b+c$. Constants may depend on $K$.

\begin{lemma}[Replacement for {\cite[Lemma~4.4]{MT}}]\label{lem:L44}
Let $\tau\in[0.45,0.4565]$ and $s_0\in[1-2\cdot10^{-4},1+O(1/\log x)]$.
\begin{enumerate}[label=(\arabic*),leftmargin=*]
\item For all $a,b\ge0$ and $c>0$: $\int_{T_0}^{\min(N_C,T)}|ABC(\frac12+it)|\,dt\ll(N_AN_BN_C)^{1/2}(\log x)^{K_1}$.
\item Suppose $c<\tau$, $a,b\in[0.087,0.4565)$ and $c\ge0.087$. Then
\[\int_{N_C}^{T}|ABC(\tfrac12+it)|\,dt\ll(N_AN_BN_C)^{1/2}(\log x)^{K_2}.\]
Moreover, the corner classes $\cQ$ are those level classes (\S\ref{sec:levels}) with $|A|\ge N_A^{1/2}x^{-0.02}$ and $|B|\ge N_B^{1/2}x^{-0.02}$.
All other classes together contribute $\ll(N_AN_BN_C)^{1/2}x^{-10^{-5}}$.
\end{enumerate}
\end{lemma}

Compared with the $0.548$ version, only the numbers change: there $\tau\le0.4521$, the lower length bound was $0.0958$ and the saving off the
corner was $x^{-10^{-4}}$. The rest of this section proves the lemma.

\subsection{Small $t$ and long $C$}\label{sec:short}

\begin{lemma}\label{lem:Csmall}
For $1\le t\le N_C$, $|C(\frac12+it)|\ll(\log x)N_C^{1/2}/t$.
\end{lemma}
\begin{proof}
Apply Lemma~\ref{lem:T411} with $\sigma=\frac12$ at $y=N_C$ and at $y=2N_C$, and subtract. This gives
$|C|\le 2(2N_C)^{1/2}/|\frac12-it|+O(N_C^{-1/2})$, and $N_C^{-1/2}\le N_C^{1/2}/t$. For $c_n=\log n$, use partial summation on the same
identity over $(N_C,y]$ with $y\le2N_C$.
\end{proof}

\begin{proof}[Proof of \Lfp(1)]
Let $U\in[T_0,\min(N_C,T)]$. The coefficients are divisor-bounded, so $\sum|\alpha_n|^2/n\ll(\log x)^{O(1)}$.
Lemma~\ref{lem:Csmall}, Cauchy--Schwarz and Lemma~\ref{lem:MV} give
\begin{align*}
\int_U^{2U}|ABC|&\ll(\log x)^{O(1)}\frac{N_C^{1/2}}{U}(U+N_A)^{1/2}(U+N_B)^{1/2}\\
&\ll(N_AN_BN_C)^{1/2}(\log x)^{O(1)}\Big(\frac1U+\frac{1}{(UN_A)^{1/2}}+\frac1{(UN_B)^{1/2}}+\frac1{(N_AN_B)^{1/2}}\Big).
\end{align*}
Each bracketed term is $\le1$. Summing over the $O(\log x)$ dyadic blocks proves (1). If $c\ge\tau$ the range in (2) is empty, so
long $C$ needs nothing more.
\end{proof}

\subsection{Pointwise caps and the transfer from $\zeta$ to $C$}

The trivial bound gives $|A(\frac12+it)|\ll N_A^{1/2}(\log x)^{O(1)}$, and likewise for $B$. For $t\ge1$, apply Lemma~\ref{lem:T3}
to $f(n)=-\frac{t}{2\pi}\log n$ on dyadic ranges, then partial summation:
\begin{equation}\label{eq:Ccap}
|C(\tfrac12+it)|\ll(\log x)\big(t^{1/6}+N_C^{1/2}t^{-1/6}\big).
\end{equation}

\begin{lemma}[Transfer lemma]\label{lem:Z}
If the $t_r\in[N_C,T]$ are $1$-separated, then $\sum_r|C(\frac12+it_r)|^4\ll T(\log x)^{12}$.
\end{lemma}

\begin{proof}
Take $c_n\equiv1$ first, and fix $t\ge N_C\ge2$. Set $y_1=\lfloor N_C\rfloor+\frac12$, $y_2\in(\lfloor N_C\rfloor,\lfloor2N_C\rfloor]+\frac12$, $\kappa=\frac12+1/\log x$ and $U=2t$.
Apply the truncated Perron formula at $y_1$ and at $y_2$ and subtract. Since $y_j\in\mathbb Z+\frac12$, the truncation error is
$\ll N_C^{1/2}(\log x)/U\ll\log x$.

Now move the contour to $\operatorname{Re}w=0$:
\begin{itemize}[leftmargin=*]
\item The pole of $\zeta(\frac12+it+w)$ at $w=\frac12-it$ contributes $\ll N_C^{1/2}/t\le1$.
\item At $w=0$, $(y_2^w-y_1^w)/w$ is regular.
\item By convexity the horizontal segments are $\ll(3t)^{1/4+o(1)}N_C^{1/2}/U\ll1$, using $t\ge N_C$.
\end{itemize}
On $\operatorname{Re}w=0$ we have $|y_2^{iv}-y_1^{iv}|/|v|\le k(v):=\min(\log3,2/|v|)$. Hence
\[
\Big|\sum_{N_C<n<y_2}n^{-1/2-it}\Big|\le\frac1{2\pi}\int_{-2t}^{2t}|\zeta(\tfrac12+i(t+v))|\,k(v)\,dv+O(\log x).
\]
The weight $\log n$ costs a factor $\ll\log x$, by partial summation in $y_2$. By H\"older, using $\int_{-2t}^{2t}k\ll\log x$,
\[
|C(\tfrac12+it)|^4\ll(\log x)^{4+3}\int_{-2t}^{2t}|\zeta(\tfrac12+i(t+v))|^4k(v)\,dv+(\log x)^{8}.
\]
For $1$-separated $t_r$, $\sum_rk(w-t_r)\ll\log x$ uniformly in $w$. Summing over $r$ and using Lemma~\ref{lem:Ing} on $[-T,3T]$ gives
$\sum_r|C(\frac12+it_r)|^4\ll(\log x)^{8}\,T\log^4T+R(\log x)^8\ll T(\log x)^{12}$.
\end{proof}

The value of $\tau$ enters this proof only through $T\le x$; so Lemma~\ref{lem:Z}, like \eqref{eq:Ccap} and \Lfp(1), holds verbatim in the
larger range of $\tau$ used here.

\subsection{Discretisation and level classes}\label{sec:levels}

Cover $[N_C,T]$ by the unit intervals $I_j=[N_C+j,N_C+j+1]$, $j\ge0$, starting at $N_C$. Let $t_j$ maximise $|ABC(\frac12+it)|$ on
$I_j\cap[N_C,T]$, so every $t_j\in[N_C,T]$. Then $\int_{N_C}^T|ABC|\le\sum_j|ABC(\frac12+it_j)|$. The points with $j$ even, and those with
$j$ odd, are each $1$-separated.

Sort the points into dyadic classes $|A(\frac12+it_j)|\in(\frac12x^{\lambda_A},x^{\lambda_A}]$, and likewise $\lambda_B$ for $B$ and
$\gamma$ for $C$. There are $O((\log x)^3)$ classes with all levels $\ge-\frac12$, plus three low classes, in which $|A|$, $|B|$ or $|C|$ is $\le x^{-1/2}$.
Write $R=x^{\rho}$ for the number of points in a class and set
\[
E:=\rho+\lambda_A+\lambda_B+\gamma-s_0/2.
\]
A class contributes $\le x^{s_0/2+E}$, so it suffices to show $E\le-\eta+o(1)$ for every class.

\emph{Clamping.} A class level may exceed its cap ($a/2$, $b/2$, or \eqref{eq:Ccap} for $C$) by $o(1)$, and $s_0$ may exceed $1$ by
$O(1/\log x)$. We clamp each class's lower level to its cap, and $c$ to the end of its range. The count bounds stay valid, since they are
non-increasing in $\lambda$. The exponent $E$ changes by $o(1)$, which the fixed margin absorbs.

\emph{Low classes.} By the caps, $\lambda_A\le a/2+o(1)$, $\lambda_B\le b/2+o(1)$ and $\gamma\le c/2+o(1)$. If, say, $\lambda_A<-\frac12$, then
$\rho\le\tau+o(1)$ gives $E\le\tau-\frac12-\frac a2+o(1)<0$.

\emph{Normalisation.} Let
\[
\cX=\{A^k,B^k\ (k\le3),\ (AB)^k\ (k\le2),\ C^k\ (k\le6),\ AC,\ BC\}.
\]
Each $X\in\cX$ has length $x^{n_X}$. Its coefficients are $\ll d_{O(1)}(n)(\log x)^{O(1)}=x^{o(1)}$, and it is supported in $O(1)$ dyadic ranges.
On a class, $|X|\ge x^{\lambda_X-o(1)}$, where $\lambda_X$ is the corresponding sum of levels (e.g.\ $\lambda_{AB}=\lambda_A+\lambda_B$).

Pigeonhole a positive proportion of the points onto one dyadic piece $\sum_{n\sim N}c_nn^{-1/2-it}$. Rescale to
$b_n=c_n(N/n)^{1/2}/\max|c_m|$, so that $|b_n|\le1$. Lemmas~\ref{lem:GM1}--\ref{lem:GM12} then apply with $V=N^{\sigma}$ and
$\sigma=\frac12+\lambda_X/n_X-o(1)$.

\emph{Box edges.} Lemma~\ref{lem:GM1} has no range hypotheses. Lemma~\ref{lem:GM12} needs $\sigma\ge7/10$ and $T^{5/6}\le N\le T$.
The certificate uses it, in the form E3k of \S\ref{sec:counts}, only on boxes where these hold at the nominal exponents, with a strict slack in $\sigma$.
\begin{itemize}[leftmargin=*]
\item \emph{The $\sigma$-edge.} E3k is used only on boxes with $\lambda_0\ge n_1/5+10^{-6}$ (exact test). On such a box $\lambda\ge\lambda_0$, $n\le n_1$
and $n_1<1$, so $\lambda/n\ge1/5+10^{-6}/n_1>1/5+10^{-6}$. Hence $\sigma=\frac12+\lambda_X/n_X-o(1)\ge7/10+10^{-6}-o(1)>7/10$ for $x$ large,
uniformly, since the boxes are finitely many and fixed. No edge argument is needed. In fact the slack used is much larger: on every box
where the certificate uses E3k, $\lambda_0-n_1/5\ge1571/2560000\approx6.1\cdot10^{-4}$ (measured by R4).
\item \emph{Why a slack.} The $0.548$ version handled $\sigma\in[7/10-o(1),7/10)$ by switching to Lemma~\ref{lem:GM1}. That works for the closed
form of Lemma~\ref{lem:GM12} but not for the $\inf_k$ display, which can be smaller than Lemma~\ref{lem:GM1} at $\sigma=7/10$: at $N=T$,
Lemma~\ref{lem:GM1} gives $N^{4/5}$ while the $k=1$ term gives $N^{0.6}+N^{0.7}+N^{0.4}+N^{5.2/7}$ (found by GA, confirmed by R4).
So that step is not used here.
\item \emph{The lower $N$-edge} is automatic. E3k is used only where $n_0\ge5\tau_{\rm hi}/6$, so every dyadic piece has $N\ge x^{n_0}\ge T^{5/6}$.
\item \emph{The upper $N$-edge.} If $N\in(T,2^{O(1)}T]$, apply Lemma~\ref{lem:GM12} with $T'=N$. Every power of $T$ in E3k has exponent in
$[0,1]$, so the bound changes by a factor $O(1)$.
\end{itemize}
In each case the exponents move by $o(1)$, which the margin absorbs.

\subsection{Count bounds and monotonicity}\label{sec:counts}

Fix $X\in\cX$ with length exponent $n$ and level $\lambda>0$, and put $\sigma=\frac12+\lambda/n$. Up to $+o(1)$, the following hold.
\begin{align*}
\text{(E1, Lemma~\ref{lem:MV})}\quad&\rho\le\max(\tau,n)-2\lambda;\\
\text{(E2, Lemma~\ref{lem:GM1})}\quad&\rho\le\max\big(n-2\lambda,\ \tfrac85n-4\lambda,\ \tau+\tfrac25n-4\lambda\big);\\
\text{(E3k, Lemma~\ref{lem:GM12}, $k\ge1$)}\quad&\rho\le\max\Big(n-2\lambda,\ \tfrac\tau2+n-4\lambda,\ \tfrac{k}{k+1}(\tau+n-6\lambda),\
\tfrac{4k}{4k+3}(2n-6\lambda)+\tfrac{2\tau}{4k+3}\Big)\\
&\qquad\text{if }\tfrac56\tau\le n\le\tau\text{ and }\lambda\ge n/5+10^{-6}\text{, i.e.\ }\sigma\ge\tfrac7{10}+10^{-6}/n;\\
\text{(E4, Lemma~\ref{lem:Z}, $X=C$)}\quad&\rho\le\tau-4\gamma;\qquad\text{and trivially}\quad\rho\le\tau.
\end{align*}
E3k is the first display of Lemma~\ref{lem:GM12} at a fixed $k$, written in exponents ($n=\log_xN$, $\tau=\log_xT$): since $\sigma n=n/2+\lambda$,
we have $(2-2\sigma)n=n-2\lambda$, $(3-4\sigma)n=n-4\lambda$, $(4-6\sigma)n=n-6\lambda$ and $(5-6\sigma)n=2n-6\lambda$.
The certificate uses $k=1,\dots,12$. The slack $10^{-6}$ in the hypothesis is the one explained under \emph{Box edges} (\S\ref{sec:levels}).
The closed form of Lemma~\ref{lem:GM12} is not among these estimates.

\emph{Monotonicity.} Every piece of E1, E2, E3k and E4 is non-decreasing in $n$ and in $\tau$ and non-increasing in $\lambda$. For E3k this is
immediate, since each piece is linear: the coefficients of $n$ are $1,1,\frac{k}{k+1},\frac{8k}{4k+3}$, those of $\tau$ are
$0,\frac12,\frac{k}{k+1},\frac{2}{4k+3}$, all $\ge0$, and those of $\lambda$ are negative. So on a box $n\in[n_0,n_1]$, $\lambda\ge\lambda_0$,
$\tau\le\tau_{\rm hi}$, each estimate is at most its value at $(n_1,\lambda_0,\tau_{\rm hi})$. Also $\sup_{\rm box}\min_e\le\min_e\sup_{\rm box}$;
in fact the certificate accepts a box on a single estimate, so no interchange of $\sup$ and $\min$ is needed (R4).

\subsection{The corner classes}\label{sec:corner}

Put $\delta_A=a/2-\lambda_A$ and $\delta_B=b/2-\lambda_B$. The corner classes are those with $\delta_A,\delta_B\le\kappa:=0.02$.

\emph{Case $\min(a,b)\ge0.276$.} Apply Lemma~\ref{lem:HM} to $A$ with $c_n=\alpha_nn^{-1/2}$, so that
$G:=\sum|c_n|^2\ll(\log x)^{O(1)}$, and take $V=\frac12N_A^{1/2}x^{-\delta_A}$.
\begin{itemize}[leftmargin=*]
\item $V^2\ge x^{a-2\kappa}/4\ge x^{0.236}/4$, which exceeds a large multiple of $GT^{1/2}\log T\le x^{0.22825+o(1)}$.
\item So $R\ll GN_AV^{-2}\log T\ll(\log x)^{O(1)}x^{2\delta_A}$. The same argument for $B$ gives $R\ll(\log x)^{O(1)}x^{2\delta_B}$.
\item Since $|C|\ll N_C^{1/2}\log x$, the class contributes
$\ll(\log x)^{O(1)}x^{s_0/2+\min(2\delta_A,2\delta_B)-\delta_A-\delta_B}\le(\log x)^{O(1)}x^{s_0/2}$.
\end{itemize}
This is a loss of a power of $\log x$ only, with no power saving. In the application it is absorbed by Lemma~\ref{lem:VK}.
This case works whenever $0.236>\tau/2$, i.e.\ for $\tau<0.472$.

\emph{Case $\min(a,b)<0.276$.}
\begin{itemize}[leftmargin=*]
\item Here $c>1-2\cdot10^{-4}-0.276-0.4565=0.2673$.
\item By \eqref{eq:Ccap}, $c/2-\gamma\ge\min(c/2-\tau/6,\,c/6)-o(1)\ge\min(0.0575,\,0.04455)-o(1)=0.04455-o(1)$.
\item Lemma~\ref{lem:MV} applied to $AB$, of length $a+b>1-2\cdot10^{-4}-\tau>\tau$, gives $\rho\le2\delta_A+2\delta_B+o(1)$.
\item Hence $E\le\delta_A+\delta_B-0.04455+o(1)\le2\kappa-0.04455+o(1)=-0.00455+o(1)$.
\end{itemize}
This case works whenever $(0.724-2\cdot10^{-4}-\tau)/6>2\kappa$ and $(0.724-2\cdot10^{-4}-\tau)/2-\tau/6>2\kappa$, i.e.\ for $\tau<0.483$.

\subsection{The computer certificate}\label{sec:cert}

All remaining classes are handled by the script\\ \texttt{numerics/certify\_exact\_k.py}\footnote{MD5 \texttt{8c4aded7ba48509ceafb3480a2be9bf3}. Invocation:\\
{\footnotesize\texttt{TLO=0.45 THI=0.4565 MARGIN=0.00001 EPSP=0.0002 VSLACK=0.000001 MAXIT=150}}\\
{\footnotesize\texttt{python3 certify\_exact\_k.py min0k\_noclosed 0.02}}: certified after 67 rounds (iterations 0--66), $331\,054$ leaves;
log \texttt{numerics/run\_min0knc\_4565.out}.}.
It is a branch-and-bound over boxes in $(a,b,\lambda_A,\lambda_B,\gamma)$, in exact rational arithmetic. It is the exact certifier
\texttt{numerics/certify\_exact.py} of report X1 (MD5 \texttt{f083c952\dots}, unchanged and imported), with the estimates E3k added.
According to R4's diff, nothing else changed: the tiling, the empty and corner tests, the caps, the tails and the acceptance test are X1's.
One documentation slip: line~8 of the script's docstring describes the E3k validity test as $\lambda_0\ge n_1/5$, as for the closed form. The
code also requires the slack, $\lambda_0\ge n_1/5+\texttt{VSLACK}$ with $\texttt{VSLACK}=10^{-6}$, which is the test stated here (R4, issue m2; the code is correct).

\emph{Setup.}
\begin{itemize}[leftmargin=*]
\item On a box, $c$ ranges over $[1-a_1-b_1-2\cdot10^{-4},\,1-a_0-b_0]$, which covers $s_0\in[1-2\cdot10^{-4},1]$.
\item $\tau\in[\tau_{\rm lo},\tau_{\rm hi}]=[9/20,913/2000]$ is treated uniformly, and every bound uses $\tau_{\rm hi}$. The range conditions of E3k are checked as
$n_1\le\tau_{\rm lo}$, $n_0\ge5\tau_{\rm hi}/6$ and $\lambda_0\ge n_1/5+10^{-6}$ (\S\ref{sec:levels}).
\item The initial boxes are $(a,b)$-cells of width $1/50$ starting at $1-2\tau_{\rm hi}=0.087$, with $a\le b$ by symmetry, together with
$\lambda_A\in[-\frac12,a_1/2]$, $\lambda_B\in[-\frac12,b_1/2]$ and $\gamma\in[-\frac12,\frac12]$.
\end{itemize}

\emph{The test on each box.}
\begin{itemize}[leftmargin=*]
\item $\bar\rho$ is the minimum, over $X\in\cX$ and over E1, E2, E3k ($k\le12$), E4 and the trivial bound, of the box suprema from \S\ref{sec:counts}. These are evaluated at the
lower level ends. A polynomial whose level lower end is $\le0$ is not used.
\item The contribution exponent uses the upper ends of the level ranges. These are capped by $a_1/2$ and $b_1/2$, and for $C$ by
\[\min(c_1/2,\max(\tau_{\rm hi}/6,c_1/3)).\]
\item The box is certified if
$E_{\rm box}:=\max(\bar\rho,0)+\bar u_A+\bar u_B+\bar u_C-(1-2\cdot10^{-4})/2\le-10^{-5}$.
\end{itemize}

\emph{Exempt boxes.} Two kinds of boxes are not tested:
\begin{itemize}[leftmargin=*]
\item empty boxes, where a level exceeds its cap or $a,b<\tau_{\rm hi}$, $0.087\le c<\tau$ fails;
\item boxes inside $\cQ$, which \S\ref{sec:corner} covers.
\end{itemize}

\emph{The run.} Floats only decide which boxes to try and in what order; every leaf is accepted only after an exact test. Uncertified boxes
are bisected along their widest weighted side.
\begin{itemize}[leftmargin=*]
\item The run ends with no uncertified box after 67 rounds (iterations 0--66). The leaves tile the domain exactly: $331\,054$ leaves, namely $807$ empty, $3\,310$ in
the corner region and $326\,937$ certified.
\item The largest certified value is $E_{\rm box}\le-18999/1884160000\approx-1.008\cdot10^{-5}$. It sits on the five-factor box,
$(a,b)\approx(0.400,0.400)$, $\gamma\approx0.0434$.
\item The low classes (some level $<-\frac12$) are certified in closed form, with $E\le-9779/60000\approx-0.163$ and $E\le-869/10000$.
\item Among the $\inf_k$ terms, $k=5$ and $k=4$ win most often. This agrees with the five-factor analysis, which uses $k=5$.
\item \emph{Independent re-verification} (R4). The certifier stores no leaves, so its soundness rests on re-running it. The referee R4 therefore wrote
\nolinkurl{rounds/round2/R4_scripts/r4_verify.py} (MD5 \texttt{471a01049eeb3af168806a92bce61ec6}). It records every accepted leaf of the
\texttt{min0k\_noclosed} run and re-checks it with independently written code that does not call the certifier's evaluation functions:
\begin{itemize}[leftmargin=*]
\item the bounds are written in $\sigma$-form, directly from the cited statements;
\item box suprema are taken by enumerating the vertices $\{n_0,n_1\}\times\{\lambda_0,\lambda_1\}\times\{\tau_{\rm lo},\tau_{\rm hi}\}$. This is valid because
each candidate is a maximum of affine functions, and it does not rely on the monotonicity claim of \S\ref{sec:counts};
\item the E3k hypotheses are checked with strict $\sigma>7/10$;
\item the empty and corner tests, the caps and the acceptance test are recomputed;
\item the exact total volume of the leaves is checked to equal that of the initial cells, so no box is lost.
\end{itemize}
Result: $331\,054$ leaves at $\tau_{\rm hi}=913/2000$ and $586\,580$ at $\tau_{\rm hi}=114173/250000=0.456692$, with $0$ failures and with
equal volumes in both runs. The maxima agree with the certifier's. Usage: run it with \texttt{PYTHONPATH=numerics} from a scratch directory, since it writes a JSON file.
\item \emph{Re-runs.} R4 re-ran the certifier for \texttt{min0k\_noclosed} and \texttt{min0k} at $913/2000$ and for \texttt{min0k\_noclosed} at
$0.456692$, and obtained bit-identical leaf counts. The writer of this note also re-ran the headline invocation from fresh copies of the three
scripts (same MD5 sums), in a separate directory. It reproduced the same leaf counts and the same maximum.
\item Rigour of the box bounds: they follow from the monotonicity in \S\ref{sec:counts}. Since there are finitely many boxes, the $o(1)$ terms are
absorbed by the fixed margin $10^{-5}$. Levels are clamped to their caps as in \S\ref{sec:levels}.
\item Scope: the run checks only the arithmetic of the box bounds, given E1--E4, the caps and the corner lemma of \S\ref{sec:corner}.
\end{itemize}

\emph{Inputs actually used.} The run is called \texttt{min0k\_noclosed}. Its inputs are MVT, Lemma~\ref{lem:GM1}, the $\inf_k$ display of
Lemma~\ref{lem:GM12} ($k\le12$) and Lemma~\ref{lem:Z}. The following are \emph{not} used: the closed form of Lemma~\ref{lem:GM12}, Huxley's
estimate, Jutila's estimate, the twelfth moment, Bourgain's bound for $\zeta$, and Vinogradov--Korobov for $C$.

\emph{Cross-checks and sanity runs} (same code and settings). None of these is used in the proof.
\begin{itemize}[leftmargin=*]
\item \texttt{min0k} adds the closed form of Lemma~\ref{lem:GM12}, with its old test $\lambda_0\ge n_1/5$; for the closed form alone, the
$\sigma$-edge step of the $0.548$ version is valid. Result: certified, $330\,937$ leaves. Adding the full input set of the earlier code as well
(\texttt{fullk}): certified, $310\,004$ leaves. Both take 67 rounds.
\item Regression: with the old input set \texttt{min0} at $\tau_{\rm hi}=0.4521$, the new code reproduces the tiling of the $0.548$ version
exactly ($98\,433$ leaves: $661$ empty, $2\,432$ corner, $95\,340$ certified).
\item In every run, the float mirror of the new code agrees with the float code of the $0.548$ version on all old estimates, with maximum
difference $0.0$. That code is \nolinkurl{numerics/certify_548.py} (MD5 \texttt{4d0afc03\dots}).
\item The sanity run of the $0.548$ version failed at $\vartheta=0.5435$, with $\max E=+0.0019$ at $(a,b,c)=(0.4,0.4,0.2)$. That run used the
closed form of Lemma~\ref{lem:GM12}. With the closed form, the five-factor box closes only for $\vartheta>0.5437644$ (\S\ref{sec:limits}). The
present certificate, with the $\inf_k$ display, supersedes it.
\item The control run at $\vartheta=0.5433<69/127$ fails (Remark~\ref{rem:543308}).
\item In the $0.548$ version, a run without Lemma~\ref{lem:Z} failed at $(\frac13,\frac13,\frac13)$ by $+0.025$. So the transfer of the fourth
moment from $\zeta$ to $C$ is essential.
\end{itemize}

This completes the proof of \Lfp: the corner classes are handled in \S\ref{sec:corner} and all others by the certificate.\qed

\section{Proof of Theorem~\ref{thm:main}}\label{sec:proof}

\emph{Parameters.} Let $\vartheta>0.5435$ and $\e_0>0$. We may assume $\vartheta\le0.55$, since for $\vartheta>0.55$ the claim is
\cite[Theorem~1.1]{MT}.
\begin{itemize}[leftmargin=*]
\item First fix $\e:=\min(\e_0,(\vartheta-0.5435)/4,10^{-4})$ and $\theta:=\vartheta-\e$. Then $\theta\in[0.5435+3\e,0.55-\e]$. These depend only on
$\vartheta$ and $\e_0$, not on $H$. (The $0.548$ version divided by $2$; dividing by $4$ keeps $b$ strictly below $0.4565$, see below.)
\item Then, as in \cite[\S4.1]{MT}, split the sum into short sums. This reduces the theorem to intervals $(y,y+H']$ with $y\ge x$ and
$H'\in[y^{\theta+\e},2y^{\theta+\e}]\subset[y^{\e},y]$, where Lemma~\ref{lem:ramare} applies. This holds also when $H>x$. We write $x,H$ for $y,H'$.
\end{itemize} Run MT's argument (\S\ref{sec:red}) with these $\theta,\e$. By Lemma~\ref{lem:ramare} the saving is $(\log x)^{-1/3+\e}$, and
$\e\le\e_0$. By Lemma~\ref{lem:perron} it suffices to prove \eqref{eq:L1} for each configuration \eqref{eq:config}. The upper limit of
integration is $T=x^{1+\e/10}/H\le x^{\tau}$ with $\tau:=1-\theta-0.9\e=1-\vartheta+0.1\e\in[0.45,0.4565)$.

\emph{Removing $P$ and $R$.} Factors with $N_i<2$ have at most three terms and are $O(1)$; drop them. Write $N_i=x^{\alpha_i}$ and $s_0=\sum\alpha_i$. Then
$s_0\in[1-\e/2-o(1),1+O(1/\log x)]$, because $R\le x^{\e/2}$ and $P_1=x^{o(1)}$. In particular $s_0\ge1-2\cdot10^{-4}$. We have:
\begin{itemize}[leftmargin=*]
\item $|P(\frac12+it)|\ll_AP_1^{1/2}(\log x)^{-A}$ on $[T_0,T]$, by Lemma~\ref{lem:VK};
\item $|R(\frac12+it)|\ll R^{1/2}(\log x)^{O(1)}$, trivially;
\item $P_1R\,x^{s_0}\asymp x$.
\end{itemize}
So it suffices to show $\int_{T_0}^T|N_1\cdots N_{2k-1}(\frac12+it)|dt\ll x^{s_0/2}(\log x)^{O(1)}$.

\emph{Type II.} Let $w_1=1-\theta$ and $w_2=\theta-\e$. Suppose some subsum $\sum_I\alpha_i\in[w_1,w_2]$. Then $X_I\ge x^{1-\theta}\ge T$ and
$X_{I^c}\ge x^{s_0-\theta+\e}\ge x^{1-\theta}\ge T$, and Lemma~\ref{lem:typeII} applies.

\emph{Otherwise.} Apply Lemma~\ref{lem:comb} (note $s_0>w_2$). This gives $I$, $r$ and $J$ with
\begin{gather*}
a:=\textstyle\sum_I\alpha_i<1-\theta=1-\vartheta+\e<0.4565,\qquad c:=\alpha_r>2\theta-1-\e\ge0.087+5\e>0.087,\\
b:=\textstyle\sum_J\alpha_j<s_0-w_2\le1-\vartheta+2\e+O(1/\log x)\le0.4565-\tfrac12(\vartheta-0.5435)+O(1/\log x).
\end{gather*}
\begin{itemize}[leftmargin=*]
\item Since $c>1/20+o(1)$ and the M\"obius-type factors have length $\le(2x)^{1/20}$, the factor $N_r$ is smooth: $r\le k-1$.
\item Put $A=\prod_IN_i$, $B=\prod_JN_j$ and $C=N_r$. These satisfy the hypotheses of \S\ref{sec:L44} with $K=40$.
\item Lemma~\ref{lem:L44}(1) covers $[T_0,\min(N_C,T)]$.
\item If $c<\tau$, Lemma~\ref{lem:L44}(2) covers $[N_C,T]$. Its hypotheses hold:
\begin{itemize}[leftmargin=*]
\item $a>w_2-c>\theta-\e-\tau=2\theta-1-0.1\e\ge0.087+5.9\e>0.087$;
\item $b>s_0-w_1-\tau\ge2\theta-1+0.4\e-o(1)\ge0.087$;
\item $a<0.4565$; $b<0.4565$ for $x$ large, by the display above; and $0.087\le c<\tau$.
\end{itemize}
\end{itemize}
In every case the integral is $\ll x^{s_0/2}(\log x)^{O(1)}$. Combined with the factor $(\log x)^{-A}$ from $P$, this gives \eqref{eq:L1}, and
Theorem~\ref{thm:main} follows. We never use Lemma~4.4 of \cite{MT}.\qed

\begin{remark}[Closer to $69/127$]\label{rem:543308}
The same proof gives Theorem~\ref{thm:main} for somewhat smaller $\vartheta$, at the price of smaller margins. We keep the headline at
$0.5435$, where the margin is $10^{-5}$. Each item below uses the same code, input set \texttt{min0k\_noclosed} (also certified with \texttt{min0k}), $\kappa=1/50$,
and the corner lemma, which holds for $\tau<0.472$.
\begin{itemize}[leftmargin=*]
\item $\vartheta>0.5434$: $\tau_{\rm hi}=2283/5000$, \texttt{EPSP}$\,=10^{-4}$, margin $10^{-5}$ and $\e:=\min(\e_0,(\vartheta-0.5434)/4,10^{-4})$.
This \texttt{EPSP} is valid, since $s_0\ge1-\e/2-o(1)>1-10^{-4}$. The lower length bound is $0.0868$. Certified after 76 rounds with $388\,691$ leaves
($388\,580$ for \texttt{min0k}). Corner case (ii): $c>0.2673$.
\item $\vartheta>0.543308=69/127+9.1\cdot10^{-7}$: $\tau_{\rm hi}=114173/250000=0.456692$, \texttt{EPSP}$\,=2\cdot10^{-7}$, margin $3\cdot10^{-7}$,
$\sigma$-slack $10^{-7}$, and $\e:=\min(\e_0,(\vartheta-0.543308)/4,10^{-7})$, so that $s_0>1-2\cdot10^{-7}$. The lower length bound is $0.086616$.
Certified after 103 rounds with $586\,580$ leaves ($828$ empty, $3\,207$ corner, $582\,545$ certified); $586\,479$ for \texttt{min0k}.
The invocation is that of the footnote in \S\ref{sec:cert} with \texttt{THI=0.456692}, \texttt{MARGIN=0.0000003}, \texttt{EPSP=0.0000002} and
\texttt{VSLACK=0.0000001}. The largest certified value is $E_{\rm box}\le-4565757/15073280000000\approx-3.03\cdot10^{-7}$, and the slack used for E3k is at least
$\lambda_0-n_1/5\ge1.38\cdot10^{-3}$ (R4).
Corner lemma: $0.236>\tau/2=0.228346$; in case (ii), $c>0.267308$ and $c/6>0.04455>2\kappa$.
\item Runs at $\tau_{\rm hi}=0.45665$, $0.45668$ and $0.45669$ (\texttt{EPSP} $2\cdot10^{-5}$, $4\cdot10^{-6}$, $10^{-6}$) are also certified.
\item \emph{Control.} At $\tau_{\rm hi}=0.4567$, i.e.\ $\vartheta=0.5433<69/127$, with \texttt{EPSP}, margin and slack all $10^{-7}$, the run is not certified.
The number of open boxes doubled every round, and the run was stopped after round 84 with $317\,127$ open boxes and float maximum $E=+10^{-5}$.
An exact pointwise evaluation at $(a,b,c)=(0.4,0.4,0.2)$, $\sigma_A=\sigma_B=0.716525$, $\gamma=0.043303$ and $\tau=0.4567$ gives $E=+3\cdot10^{-6}>0$.
So the model itself fails there, and the certificate is not vacuous.
\item $69/127$ itself cannot be reached this way. At the binding point the slack is about $58/127-\tau_{\rm hi}$, so the margins would have to tend to $0$.
\end{itemize}
These extensions have the same status as Theorem~\ref{thm:main}. The referee R4 checked them as well: the corner numbers and the $\e\le10^{-7}$ bookkeeping
at $0.456692$, a re-run, and the leaf-by-leaf re-verification of \S\ref{sec:cert}. R4 recommends keeping them as a remark rather than in the theorem.
\end{remark}

\section{Limits of the method}\label{sec:limits}

The headline exponent $0.5435$ lies $1.9\cdot10^{-4}$ above $69/127=0.5433071\ldots$. Remark~\ref{rem:543308} reduces the gap to
$9.1\cdot10^{-7}$. Below $69/127$ the method itself fails, for the reason explained next.

\emph{Five smooth factors.} The configuration $N_1=\dots=N_5=x^{1/5}$ with smooth coefficients has no subsum in $[1-\theta,\theta]$
for $\theta<3/5$. So it must be handled as $A=N_1N_2$, $B=N_3N_4$, $C=N_5$, i.e.\ $(a,b,c)=(0.4,0.4,0.2)$. The maximum of the certificate and the
failing control run both sit at this point. Internal report P2 treats this configuration; an internal referee (R1) judged it correct with
fixable issues. Report PB1 (phase~II) recomputed the binding point in exact rational arithmetic.
\begin{itemize}[leftmargin=*]
\item The inputs are: MVT, Lemma~\ref{lem:T411}, the cap \eqref{eq:Ccap}, the fourth-moment bound for the fifth factor, and the $\inf_k$ form of Lemma~\ref{lem:GM12} with $k=5$.
With these inputs the five-factor case closes exactly when $\vartheta>69/127=0.5433071\ldots$, i.e.\ $\tau<58/127$.
\item At $\vartheta=69/127$ the pair products $N_iN_j$ have length $x^{2/5}\approx T^{0.88}$, the binding level is $\sigma=91/127$, each
$|N_i|$ has size $x^{11/254}$, and the number of points must stay below $x^{36/127}$. Two bounds equal $36/127$ exactly: MVT for a pair,
equivalently Lemma~\ref{lem:Z} for $N_5$; and the part $N^{(5-6\sigma)20/23}T^{2/23}$ of the $k=5$ term. The term $T^{1/2}N^{3-4\sigma}$ gives $179/635$,
only $1/635\approx0.0016$ lower. (The writer re-derived these values by hand from E1, E3k and E4.)
\item With the closed form of Lemma~\ref{lem:GM12} instead, the threshold is $\vartheta>0.5437644$, where $\tau=1-\vartheta$ is the smaller root of
$30\tau^2-32.8\tau+8.72=0$. This is why the $0.548$ version, which used the closed form, failed at $0.5435$ in its sanity run.
\end{itemize}

\emph{Independent certification.} Internal report P4 used code written independently of P3's. It certified the generic level-set model, with $A,B$ arbitrary
and $C$ smooth. Its inputs were MVT, Huxley, Lemma~\ref{lem:GM1}, the $\inf_k$ form of Lemma~\ref{lem:GM12} ($k\le8$), Jutila ($k=3$),
the fourth and twelfth \cite{HB} moments of $\zeta$, and Bourgain's bound \cite{Bo}. The results are in Table~\ref{tab:p4}.
\begin{itemize}[leftmargin=*]
\item $\vartheta=0.5432$ was not certified within the box budget. This agrees with the exact certificate of \S\ref{sec:cert}, which passes at $0.5434$ and fails at $0.5433$.
\item With pre-Guth--Maynard inputs only (MVT, Huxley, Jutila), the same code certifies $0.552$ and fails at $0.548$. This is consistent with the exponent
$0.55$ of \cite{MT}.
\item P4 certifies only the level-set model. Its reduction to the top $t$-scale and its corner argument, via Huxley in Hal\'asz form, were argued by hand and
classed as routine.
\end{itemize}

\begin{table}[h]
\centering
\begin{tabular}{@{}lcccccc@{}}
\toprule
$\vartheta$ & $0.548$ & $0.545$ & $0.544$ & $0.5438$ & $0.5435$ & $0.5434$\\
\midrule
certified margin $\eta$ & $10^{-3}$ & $10^{-3}$ & $3\cdot10^{-4}$ & $3\cdot10^{-4}$ & $10^{-4}$ & $3\cdot10^{-5}$\\
boxes & 12\,164 & 30\,298 & 33\,477 & 39\,149 & 43\,639 & 46\,242\\
\bottomrule
\end{tabular}
\caption{Independent certification of the level-set model (report P4, \texttt{numerics/p4\_check.py}).}\label{tab:p4}
\end{table}

\emph{What blocks $69/127$: findings of phase~II.} The statements below are model computations and structural observations from internal
reports (PB1, PB2, PC1, SC5, DR4, S4). Unless marked, they were not re-derived by the writer, and they are not theorems.
\begin{enumerate}[label=(\alph*),leftmargin=*]
\item \emph{Reduction to one zeta segment} (PB1, and independently the control agent PC1, which had no access to the project's idea registry).
Take all five smooth factors in the same dyadic block. This is a legitimate configuration of \eqref{eq:config}. Then $N_1=\dots=N_5=S_N$, a
single segment $S_N(s)=\sum_{n\sim N}n^{-s}$ with $N\asymp x^{1/5}$, and $T=N^{5(1-\vartheta)}$. What is needed is
\[
\int_{T_0}^{T}|S_N(\tfrac12+it)|^5\,dt\ll N^{5/2}(\log N)^{B}.
\]
So below $69/127$ the problem contains a fifth-moment problem for one segment of $\zeta$. In large-values language (notation of the ANTEDB \cite{TTY}), PB1
reports that going below $69/127$ needs $\mathrm{LV}_\zeta(\sigma,2.3)<5-5\sigma$ for $\sigma\in(0.70,0.72]$. The record values it lists are
$1.48$, $1.46$ and $1.435$ at $\sigma=0.705$, $0.71$ and $0.715$, against targets $1.475$, $1.45$ and $1.425$. So one would need to beat the record
by $0.005$--$0.010$. This is the obstacle O4sym of the project database.
\item \emph{The joint route is refuted} (PB1, SC5). The 0.548 version suggested a first-moment bound
$\sum_{t\in S_A(\sigma)}|C(\frac12+it)|\lessapprox|S_A(\sigma)|$ for $C$ on the large-value set $S_A(\sigma)$ of $A$; this would close the
five-factor family for every $\vartheta>1/2$. In the equal-segment configuration, however, $A=N_1N_2=C^2$. So $A$ is large exactly where $C$ is
large, and the bound would force a pointwise bound of about $T^{0.065}$ for a zeta segment, whereas the best known is about $T^{0.155}$ (cf.\
\cite{Bo}). This also explains why an experiment at $x=10^{10}$ found $|C|$ systematically larger on $S_A(\sigma)$. The observation came from
the project's set E9 (joint large values of $A$ and $C$).
\item \emph{Additive energy inside the Guth--Maynard method} (PB2). In the proof of \cite[Proposition~12.1]{GM}, the $k$-dependent terms arise
from one sum (called $S_2$ in report PB2). Even if an additive-energy input removed these terms entirely, the model frontier would be
$19/35=0.542857$. It is a quadruple point at $\sigma=5/7$, where MVT, the generic term $N^{5-6\sigma}$, the term $T^{1/2}N^{3-4\sigma}$
and the requirement meet; $19/35$ solves $5\vartheta-2=1-5(1-\vartheta)/8$. So this route gains at most $69/127-19/35\approx0.00045$. Below
$19/35$ no value of the energy helps, and the energy is not binding at $69/127$.
\item \emph{Range of $\sigma$} (obstacle O2a, from the $0.548$ version). Below $\approx0.5433$ the relevant range of $\sigma$ extends below $7/10$,
outside Proposition~12.1. In a sampled version of the model, Montgomery's large values conjecture on $\sigma\in[0.70,0.73]$ alone gains only
about $0.004$. The full large values conjecture would push the model to $\vartheta\to1/2$, apart from the corners.
\item \emph{Priced targets} (entry N7 of the project database; reports PB2, DR4, S4 and PC1). Table~\ref{tab:prices} lists what some
hypothetical inputs would give in the model. These are prices, not results.
\end{enumerate}

\begin{table}[h]
\centering\small
\begin{tabular}{@{}lc@{}}
\toprule
hypothetical input (exponent model) & effect on the threshold\\
\midrule
none (this note) & $69/127=0.54331$\\
$k$-dependent terms of GM's $S_2$ removed (best case of an energy input) & $19/35=0.54286$\\
a large values bound at the binding level beating the fourth-moment level by $T^{-0.002}$ & $0.5424$\\
\quad the same with $T^{-0.005}$ & $0.5410$\\
\quad the same with $T^{-0.01}$ & $0.5389$\\
MVT for pairs of segments improved by $T^{-0.0117}$ & $0.5382$\\
sixth moment of $\zeta$ on $\{|\zeta|\ge T^{\beta}\}$ bounded by $T^{1+o(1)}$ down to $\beta=0.09$ & gain $0.0033$\\
full sixth moment $\int_0^T|\zeta(\frac12+it)|^6dt\ll T^{1+o(1)}$ & gain $0.0165$\\
\bottomrule
\end{tabular}
\caption{Model prices for going below $69/127$ (internal reports; not theorems). For the restricted sixth moment, the bound is known with
$\beta=11/72$ \cite{TTY}.}\label{tab:prices}
\end{table}

\subsection{Below $69/127$: a refereed reduction}\label{sec:sub542}
The following result (SUB542) comes from report GA2. It was refereed by R5 with verdict \emph{accept with corrections} (one definitional
correction, M1, and two minor ones, m2 and m3). We state it in R5's corrected form. It is a reduction, not an improvement of
Theorem~\ref{thm:main}: the residual family contains the five-factor configuration that blocks $69/127$.

\begin{proposition}[SUB542]\label{prop:sub542}
Let $\theta=0.542$, $\tau_{\rm hi}=0.458$ and $\mathrm{LOW}=1-2\tau_{\rm hi}=0.084$, and put
$\mathrm{NE}[R]:=\{(a,b):\max(|a-0.4|,|b-0.4|)\le R\}$.
\begin{enumerate}[label=(\roman*),leftmargin=*]
\item The conclusion of \Lfp(2), with $\tau\in[0.45,0.458]$, $s_0\ge1-10^{-4}$, $a,b\in(0.084,0.458)$ and $0.084<c<\tau$, holds for every
triple with $(a,b)\notin\mathrm{NE}[0.0075]$. The proof is an exact certificate: the code of \S\ref{sec:cert}, changed only so that boxes
contained in the closed box $\mathrm{NE}[0.0075]$ are set aside as ``excluded'' and never counted as certified
(\nolinkurl{numerics/certify_exact_k_sub.py}, MD5 \texttt{00767d5a\dots}; \texttt{EPSP}$\,=10^{-4}$, margin $10^{-5}$, slack $10^{-6}$, $\kappa=1/50$).
With \texttt{min0k} it uses $425\,506$ leaves ($698$ empty, $2\,835$ corner, $418\,542$ certified, $3\,431$ excluded) in 65 rounds, with largest
certified value $-13/1280000$. Without the closed form (\texttt{min0k\_noclosed}) it uses $425\,560$ leaves. R5 re-ran both runs; the first is
bit-identical to GA2's log.
\item Consequently, for $\vartheta\in(0.542,0.55]$ and $\e:=\min(\e_0,(\vartheta-0.542)/4,10^{-4})$, the argument of \S\ref{sec:proof} closes every
configuration \eqref{eq:config} except those in the residual set $\mathrm{Res}$. A configuration is in $\mathrm{Res}$ if it has no subsum in
$[1-\theta,\theta-\e]$, every smooth factor has length $<\tau$, and every admissible split $(a,b,c)$ satisfies $(a,b)\in\mathrm{NE}[0.0075]$.
Configurations with a smooth factor of length $\ge\tau$ are closed by \Lfp(1); this is R5's correction M1. For $\alpha\in\mathrm{Res}$ the MT split of
\S\ref{sec:proof} exists and is admissible. Hence Theorem~1.1 of \cite{MT} for every $\vartheta>0.542$ reduces to \Lfp(2) on the near-equal family
$\mathrm{NE}[0.0075]$.
\end{enumerate}
\end{proposition}

\emph{The residual family} (R5). A configuration in $\mathrm{Res}$ has a smooth factor $C$ of length $c\in[0.185-\e/2-o(1),\,0.215+o(1)]$
and two groups $A,B$ of total lengths $0.4\pm0.0075$. No factor in $A$ or $B$ has length in $(0.015,\theta-\e-b]$, an interval which contains
$(0.015,0.1345]$. For exactly five smooth factors of lengths $0.2+\delta_i$, the configuration is open if and only if $|\delta_i+\delta_j|\le0.0075$ for
all $i\ne j$. This forces $|\delta_i|\le4R/3+O(\e)=0.01$. R5 proved this criterion in general and checked it by brute force on $40\,000$ random
configurations, with no mismatch.

GA2 also certified the analogous statements with $R=0.0125$ at $\theta=0.541$ ($655\,573$ leaves) and with $R=0.02$ at $\theta=0.540$
($532\,758$ leaves). These two runs were not re-run by the referee.

\emph{What the remaining case needs.} The following are model-level analyses from internal reports (entries O4sym, PC3-EQ, PB4-F and BB-MAP
of the project database), not theorems.
\begin{itemize}[leftmargin=*]
\item The residual case is the equal-segment problem of item~(a) above. The one live target is to beat the MVT (fourth-moment) count by a power
$T^{-\delta}$ for products of two smooth segments of length $x^{1/5}$ at $\sigma\in[0.709,0.717]$. A restricted moment of order about $4.03$
($4.08$ at $\theta=0.540$) would suffice; the sixth moment is not needed.
\item For uniform slices, the control agent PC3 finds this equivalent, up to $T^{\e}$, to improving Ingham's large values bound for $\zeta$ itself.
That bound gives $\ll T^{1-4v+o(1)}$ points with $|\zeta(\frac12+it)|\ge T^{v}$, and the relevant heights are $T^{v}$ with $v\in[0.0909,0.0952]$, i.e.\
$T^{0.09\text{--}0.095}$. No such improvement is known to us; PC3 states this from memory, unchecked.
\item A multi-scale, Euler-product coherence statement would suffice (PB3-ITE). It says: if $|S_N(\frac12+it)|\ge V=T^{v}$, then some segment
$S_{N'}$ with $N'\le T^{1/3}$ has $|S_{N'}(\frac12+it')|\ge V^{1-\delta_I}$ at a point with $|t-t'|\le2$. With $k=4$ and $\delta_I\le0.496$ it would give
$\vartheta^*\approx0.540$ (PB4).
\item The statement is strongly supported by numerics: the median of $1-\delta_I$ is $0.95$--$0.96$ for $N=100$--$316$ and $T\le5.6\cdot10^5$, though
only at the scale of typical fluctuations. It follows from the Lindel\"of hypothesis, but no unconditional proof is known. The obvious route
through the approximate functional equation is circular (PB3).
\end{itemize}

\section{Provenance and status}\label{sec:prov}

This note was produced by an automated multi-agent proof search run by Claude (Anthropic). It followed a strategy
co-authored by Siddharth Iyer. It is the phase~II version of a note that proved the same statement for $\vartheta>0.548$.

\emph{Sources.} The mathematics comes from internal reports, written and checked by separate agents. Phase~I: the audit of \cite{MT} (S2) and
its re-check (S3); the five-factor lemma (P2, refereed by R1); the general \Lfp\ and the assembly (P3, refereed by R2); the exact-arithmetic
certificate (X1); and the independent certification (P4). Phase~II: report GA (the $\inf_k$ estimate, the certificate
\texttt{certify\_exact\_k.py}, the box-edge correction, the corner numbers at the new $\tau$, and the parameter choice); and the limit analyses PB1,
PB2, PC1 (control), SC5, DR4 and S4. Later phase~II rounds contributed the following:
\begin{itemize}[leftmargin=*]
\item GA2, the sub-family certificate of \S\ref{sec:sub542}, refereed by R5;
\item PB3, PB4, PC3 (control) and BB, the analyses quoted at the end of \S\ref{sec:limits}.
\end{itemize}
The writer of this note (W2) collated these reports. The writer also re-derived by hand: the exponent form
of the $\inf_k$ display and its monotonicity; the $\sigma$-slack of the box edges and the counterexample to the old edge argument; the corner
numbers; the parameter inequalities of \S\ref{sec:proof}; and the values at the five-factor binding point. The writer re-ran the headline
certificate and obtained the same result. Statements that the reports do not support were omitted or are marked above.

\emph{Refereeing.} The $0.548$ version: an independent agent (R2), who did not write the proof, refereed the proofs of \Lfp\ and of the
theorem. R2 re-ran the float certificate and checked Lemmas~\ref{lem:GM1} and~\ref{lem:GM12} against the arXiv text of \cite{GM}. Its
verdict was \emph{correct with minor fixable issues}. There were eight issues, all addressed: clamping to the caps; a numerical slip; fixing
$\theta$ before splitting $H$; the move of the Vinogradov--Korobov bound to $\operatorname{Re}s=\frac12$; log powers in Lemma~\ref{lem:Z};
theorem numbers; reproducibility of a float run; and an optional exact rerun, done by X1. A fresh-eyes referee (R3) then read that note; its
minor issues were applied. The phase~II changes are:
\begin{enumerate}[label=(\arabic*),leftmargin=*]
\item the $\inf_k$ display of Lemma~\ref{lem:GM12} as an input (E3k);
\item the box-edge argument: the old $\sigma$-edge step is invalid for E3k and is replaced by the strict slack;
\item the new certificate (\S\ref{sec:cert}) and Remark~\ref{rem:543308};
\item the corner numbers and the range of \Lfp;
\item the parameter choice $\e\le(\vartheta-0.5435)/4$.
\end{enumerate}
They were refereed by an independent agent (R4), who did not write them.

\smallskip
\noindent\textbf{Referee R4.} Verdict: \emph{accept}, with no blocking or major issues. R4 found that \Lfp\ holds for
$\tau\in[0.45,0.4565]$ and that Theorem~\ref{thm:main} follows; it also accepted Remark~\ref{rem:543308} and recommended keeping it as a remark.
In detail, R4:
\begin{itemize}[leftmargin=*]
\item read the $\inf_k$ display of \cite[Proposition~12.1]{GM} in the arXiv text and re-derived its exponent form;
\item checked GA's code diff against X1 (four changes, all minimal and exact);
\item confirmed that the old $\sigma$-edge step fails for E3k and that the strict slack is necessary and sufficient;
\item recomputed the corner numbers at $\tau=0.4565$ and $0.456692$, and checked that Lemma~\ref{lem:Z} is independent of $\tau$;
\item re-derived every inequality of the parameter choice;
\item re-ran the certificate, with bit-identical leaf counts;
\item re-verified every accepted leaf with the independent script \nolinkurl{r4_verify.py}: $331\,054$ leaves at $\tau_{\rm hi}=913/2000$ and $586\,580$ at
$\tau_{\rm hi}=0.456692$, with $0$ failures (\S\ref{sec:cert}).
\end{itemize}
R4 raised four minor issues, all addressed in this version:
\begin{itemize}[leftmargin=*]
\item m1: state E3k with the strict slack, cite \texttt{min0k\_noclosed} as the primary certificate, and drop the closed form and the old edge step;
\item m2: the script's docstring omits the slack (noted in \S\ref{sec:cert});
\item m3: editorial consistency of the sanity runs, the old $0.544$ remark (now removed), \S\ref{sec:limits} and the certificate footnote;
\item m4: cite the leaf-by-leaf re-verifier for reproducibility.
\end{itemize}

\smallskip
The note has \emph{not} been refereed by humans. The main residual risks are:
\begin{enumerate}[label=(\roman*),leftmargin=*]
\item The audit of the MT text (S2) relies on quotes obtained through a summarising web tool. An independent re-check (S3) agrees, and also used such a tool. In particular,
the claim that \cite[Lemma~4.2]{MT} and the proof of \cite[Lemma~4.3]{MT} do not use $\theta>0.55$ rests on that reading, together with our own
re-derivation. R2 lists this as the one inherited caveat.
\item Classical results are cited as standard, without theorem numbers (Remark~\ref{rem:sources}).
\item The new certificate was written by one agent (GA). Its code diff was reviewed by R4, and R4's independent script re-verified its leaves.
Like the earlier versions, it checks only the arithmetic of the box bounds, given the analytic inputs.
\item The result now depends on the $\inf_k$ display of \cite[Proposition~12.1]{GM}, not only on its closed form; the closed form alone stops at
$0.5437644$.
\end{enumerate}
We have not searched the literature exhaustively for prior applications of \cite{GM} to this problem.

{\small
\begin{thebibliography}{99}
\bibitem[Bo]{Bo} J.~Bourgain, \emph{Decoupling, exponential sums and the Riemann zeta function}, J. Amer. Math. Soc. \textbf{30} (2017), 205--224.
\bibitem[Ga]{Ga} P.~X. Gallagher, \emph{A large sieve density estimate near $\sigma=1$}, Invent. Math. \textbf{11} (1970), 329--339.
\bibitem[GK]{GK} S.~W. Graham and G.~Kolesnik, \emph{Van der Corput's Method of Exponential Sums}, London Math. Soc. Lecture Note
Ser. \textbf{126}, Cambridge Univ. Press, 1991.
\bibitem[GM]{GM} L.~Guth and J.~Maynard, \emph{New large value estimates for Dirichlet polynomials}, Ann. of Math. \textbf{203}
(2026), 623--675 (journal data as listed in our problem statement); version used: arXiv:2405.20552v2.
\bibitem[Ha]{Ha} G.~Harman, \emph{Prime-Detecting Sieves}, London Math. Soc. Monographs \textbf{33}, Princeton Univ. Press, 2007.
\bibitem[HB]{HB} D.~R. Heath-Brown, \emph{The twelfth power moment of the Riemann zeta-function}, Quart. J. Math. Oxford (2) \textbf{29} (1978), 443--462.
\bibitem[HBI]{HBI} D.~R. Heath-Brown and H.~Iwaniec, \emph{On the difference between consecutive primes}, Invent. Math.
\textbf{55} (1979), 49--69.
\bibitem[Hu]{Hu} M.~N. Huxley, \emph{On the difference between consecutive primes}, Invent. Math. \textbf{15} (1972), 164--170.
\bibitem[In]{In} A.~E. Ingham, \emph{Mean-value theorems in the theory of the Riemann zeta-function}, Proc. London Math. Soc. (2)
\textbf{27} (1926), 273--300.
\bibitem[Iv]{Iv} A.~Ivi\'c, \emph{The Riemann Zeta-Function}, Wiley, New York, 1985.
\bibitem[Ju]{Ju} M.~Jutila, \emph{Zero-density estimates for $L$-functions}, Acta Arith. \textbf{32} (1977), 55--62.
\bibitem[MT]{MT} K.~Matom\"aki and J.~Ter\"av\"ainen, \emph{On the M\"obius function in all short intervals}, J. Eur. Math. Soc.
\textbf{25} (2023), 1207--1225; arXiv:1911.09076.
\bibitem[Mo]{Mo} H.~L. Montgomery, \emph{Topics in Multiplicative Number Theory}, Lecture Notes in Math. \textbf{227}, Springer, 1971.
\bibitem[MV]{MV} H.~L. Montgomery and R.~C. Vaughan, \emph{Hilbert's inequality}, J. London Math. Soc. (2) \textbf{8} (1974), 73--82.
\bibitem[Ti]{Ti} E.~C. Titchmarsh, \emph{The Theory of the Riemann Zeta-Function}, 2nd ed., revised by D.~R. Heath-Brown, Oxford
Univ. Press, 1986.
\bibitem[TTY]{TTY} T.~Tao, T.~Trudgian and A.~Yang, \emph{Analytic Number Theory Exponent Database (ANTEDB)}, online,
\url{https://teorth.github.io/expdb/}.
\end{thebibliography}
}

\end{document}
